%%%%%
%
\section{Rotation horizontale, du rep�re g�ographique local vers un rep�re mobile}

\itemstyle{Op\'erateur~:} {\sf Geo2mobile\_2D}
\index{Geo2mobile\_2D}

%%%%%OBJET%%%%
\itemstyle{Objet~:}
   Applique une rotation � des coordonn�es cart�siennes 
\textbf{horizontales}, pour passer d'un rep�re g�ographique � un 
rep�re port� par un mobile orient� vers un certain cap.

%%%%%ARGUMENTS%%%%
\bigskip
%%%%%ENTREES%%%%
\begin{tabarg}
\multicolumn{3}{l}{\titretab {Entr\'ees}}\\
\hline
$X_e$ & Vecteur, {\tt float} & Composante sur l'axe Ouest-Est, positif vers l'Est (unit� quelconque)\\
$X_n$ & Vecteur, {\tt float} & Composante sur l'axe Sud-Nord, positif vers le Nord (m�me unit� que $X_e$)\\
$\psi$ & Vecteur, {\tt float} & Cap g�ographique (degree)\\
\end{tabarg}

\medskip
%%%%%SORTIES%%%%
\begin{tabarg}
\multicolumn{3}{l}{\titretab {Sorties}}\\
\hline
$Y_t$ & Vecteur, {\tt float} & Composante sur l'axe transverse, positif vers la droite (m�me unit� que $X_e$ et $X_n$)\\
$Y_l$ & Vecteur, {\tt float} & Composante sur l'axe longitudinal, positif vers l'avant (m�me unit� que $X_e$ et $X_n$)\\
\end{tabarg}


%%%%%FORMULE%%%%
\itemstyle{Formule}

%%Geo2mobile
\[
  \left( \begin{array}{c} Y_t \\ Y_l \end{array} \right)
  = 
  \left( \begin{array}{cc} \cos\psi & -\sin\psi \\ \sin\psi & \cos\psi \end{array} \right)
  \left( \begin{array}{c} X_e \\ X_n \end{array} \right)
\]



%%%%%COMMENTAIRE%%%%
\itemstyle{Commentaire~:}
   Il y a une l�g�re diff�rence avec une matrice de 
rotation classique car le cap est mesur� dans le sens
r�trograde � partir de l'axe des Y.

%%%%%AUTEUR%%%%
\itemstyle{Auteur~:}
CNRM/GMEI/TRAMM

%%%%%REFERENCE%%%%
\itemstyle{R\'ef\'erences~:}

